% BEGIN INCLUDED SGA7I_English_current_buildable_source_20260823/components/169_expose_IX_section_11_6_opening_through_lemma_11_6_13.tex
% Expose IX, section 11.6 opening through Lemma 11.6.13 (statement).
% French transcription lines 3192--3299; authority physical pp. 468--471
% (scan indices 467--470; source folios 456--459). Physical p. 471 was
% checked through the exact proof boundary. All pages were compared in four
% overlapping 1100-dpi bands. Checks, primes, subscripts, hooks, underlines,
% quotient bars, and formula punctuation are clear at that resolution.

\medskip
\noindent
\SGARefTarget{sga7.sga7i.eix.n11-6}{11.6}. 为定义同构 (\SGARefLink{sga7.sga7i.eix.n11-5-1}{11.5.1}), 可限于 $S$ 严格局部的情形 (\SGARefLink{sga7.sga7i.eix.n11-0}{11.0}); 此时 $\underline M$ 和 $\underline M'$ 为常层, 值为

\begin{equation}\tag{\SGARefTarget{sga7.sga7i.eix.n11-6-1}{11.6.1}}
M=D(T'_o),\qquad M'=D(T_o),
\end{equation}
其中 $T_o$ 和 $T'_o$ 分别为 $A_o$ 和 $A'_o$ 的极大环面. 同样将 $\Phi_o$ 等同于通常的有限群后, 必须在其分量 $\Phi_o(\ell)$ 与

\begin{equation*}
\Phi(\ell)=\operatorname{coker}\bigl(
 \widetilde u_\ell:M_\ell\longrightarrow\check M'_\ell\bigr).
\end{equation*}
之间定义典范同构. 实际上, 我们将直接定义典范同构

\begin{equation}\tag{\SGARefTarget{sga7.sga7i.eix.n11-6-2}{11.6.2}}
\Phi_o(\ell)\xrightarrow{\ \sim\ }
 \ker\bigl(\widetilde u_\ell\otimes D_\ell:
 M_\ell\otimes D_\ell\longrightarrow
 \check M'_\ell\otimes D_\ell\bigr).
\end{equation}
由于 $\Phi_o(\ell)$ 有限且 $M_\ell$ 与 $\check M'_\ell$ 同秩, 首先将得出 $\widetilde u_\ell:M_\ell\to\check M'_\ell$ 为同源. 随后可知 (\SGARefLink{sga7.sga7i.eix.n11-6-2}{11.6.2}) 的右端典范同构于 $\Phi(\ell)$:

\begin{equation}\tag{\SGARefTarget{sga7.sga7i.eix.n11-6-3}{11.6.3}}
\Phi(\ell)=\operatorname{coker}\widetilde u_\ell
 \xrightarrow{\ \sim\ }
 \ker(\widetilde u_\ell\otimes D_\ell).
\end{equation}
事实上, 利用 $\widetilde u_\ell$ 将 $M_\ell$ 等同于 $\check M'_\ell$ 的有限指数子模后, 可将 $\widetilde u_\ell\otimes D_\ell$ 的源和目标写成\footnote{此处印本将 $V/M_\ell$ 标为 $\Phi(\ell)$, 将 $V/\check M'_\ell$ 标为 $\Phi'(\ell)$. 实际上, 这两个商是下述映射的源和目标, 其核为有限余核 $\check M'_\ell/M_\ell$; 这里不重复这些相互矛盾的标记.}

\begin{equation*}
M_\ell\otimes D_\ell=V/M_\ell,\qquad
\check M'_\ell\otimes D_\ell=V/\check M'_\ell,
\end{equation*}
其中

\begin{equation*}
V=M_\ell\otimes\mathbb Q_\ell
 \xrightarrow{\ \sim\ }\check M'_\ell\otimes\mathbb Q_\ell.
\end{equation*}
将 (\SGARefLink{sga7.sga7i.eix.n11-6-2}{11.6.2}) 与 (\SGARefLink{sga7.sga7i.eix.n11-6-3}{11.6.3}) 复合, 立即得到要定义的同构 (\SGARefLink{sga7.sga7i.eix.n11-5-1}{11.5.1}), 故一切归结为定义典范同构 (\SGARefLink{sga7.sga7i.eix.n11-6-2}{11.6.2}).

将 $M_\ell\otimes D_\ell$ (相应地 $\check M'_\ell\otimes D_\ell$) 写成 $M_n$ (相应地 $\check M'_n$) 的归纳极限, 其中 $n$ 遍历 $\ell$ 的幂, 则这等价于定义典范同构

\begin{equation}\tag{\SGARefTarget{sga7.sga7i.eix.n11-6-4}{11.6.4}}
{}_n\Phi_o\xrightarrow{\ \sim\ }
 \ker\bigl(\widetilde u_n:M_n\longrightarrow\check M'_n\bigr),
 \qquad n=\ell^{r+1},
\end{equation}
使其在 $n$ 变化时与显然的过渡态射相容. 这里略微滥用记号, 以 $\widetilde u_n$ 表示将 $\widetilde u_\ell:M_\ell\to\check M'_\ell$ 模 $n$ 约化得到的同态.

为此考虑 $S$ 上的群概形 ${}_nA$. 由半稳定约化的假设, 它在 $S$ 上平坦且拟有限 (\SGARefLink{sga7.sga7i.eix.n2-2-1}{2.2.1}). 有典范滤过

\begin{equation}\tag{\SGARefTarget{sga7.sga7i.eix.n11-6-5}{11.6.5}}
{}_nA\supset({}_nA^o)^f\supset({}_nA^o)^t\supset0,
\end{equation}
其中上标 $f$ (相应地 $t$) 表示相应群概形的“不变部分” (\SGARefLink{sga7.sga7i.eix.n2-2-3}{2.2.3}) (相应地“环面部分” (\SGARefLink{sga7.sga7i.eix.n5-1}{5.1})). 在 \SGARefLink{sga7.sga7i.eix.n5-1}{5.1} 中定义环面部分时, 曾假设 $S$ 完备; 这里可以如此假设 (\SGARefLink{sga7.sga7i.eix.n11-0}{11.0}). (此外, \SGARefLink{sga7.sga7i.eix.n6-1}{6.1} 中已经看到, 若有需要, 这一构造可推广到 Hensel 情形.) 置

\begin{equation}\tag{\SGARefTarget{sga7.sga7i.eix.n11-6-6}{11.6.6}}
{}_n\Psi={}_nA/({}_nA^o)^f.
\end{equation}
这是 $S$ 上的平坦分离拟有限群概形 (SGA 3 $\mathrm{VI}_A$ 3.2 及 $\mathrm{VI}_B$ \SGARefLink{sga7.sga7i.eix.n9-2}{9.2}), 其泛纤维为

\begin{equation*}
{}_n\Psi_K=({}_nA_K)/({}_nA_K)^f,
\end{equation*}
其中按定义, $({}_nA_K)^f=({}_nA^o)^f_K$ 表示 Barsotti--Tate pro 群 $T_\ell(A_K)^f$ 的第 $r$ 个组成部分 $(n=\ell^{r+1})$. 因而 (\SGARefLink{sga7.sga7i.eix.n5-5-8}{5.5.8}) 中的第一个关系给出典范同构

\begin{equation}\tag{\SGARefTarget{sga7.sga7i.eix.n11-6-7}{11.6.7}}
{}_n\Psi_K\xrightarrow{\ \sim\ }\underline M_{nK},
\qquad
\underline M_n\mathrel{\mathop:}=\underline M
 \otimes_{\mathbb Z}(\mathbb Z/n\mathbb Z).
\end{equation}
事实上, 有:

\medskip
\noindent
\underline{引理} \SGARefTarget{sga7.sga7i.eix.n11-6-8}{11.6.8}. \underline{(\SGARefLink{sga7.sga7i.eix.n11-6-6}{11.6.6}) 所定义的 $S$ 上群概形 ${}_n\Psi$ 在 $S$ 上\emph{平展}, 且同构 (\SGARefLink{sga7.sga7i.eix.n11-6-7}{11.6.7}) 唯一延拓为态射}

\begin{equation}\tag{\SGARefTarget{sga7.sga7i.eix.n11-6-8-1}{11.6.8.1}}
{}_n\Psi\lhook\joinrel\longrightarrow\underline M_n
\end{equation}
\underline{它是开浸入.}

为证明 ${}_n\Psi$ 平展, 由 (\SGARefLink{sga7.sga7i.eix.n11-6-7}{11.6.7}), 只需在其特殊纤维的各点处验证这一点, 或等价地验证此纤维平展. 这是显然的, 因为所论纤维为

\begin{equation*}
{}_nA_o/{}_nA_o^o\subset A_o/A_o^o=\Phi_o.
\end{equation*}
可知 ${}_n\Psi$ 正规; 又由于 $\underline M_n$ 在 $S$ 上有限, 延拓 (\SGARefLink{sga7.sga7i.eix.n11-6-8-1}{11.6.8.1}) 的存在立即得出. 这是 $S$ 上平展群概形之间的态射, 故其核在 $S$ 上平展. 由于 ${}_n\Psi$ 从而此核分离, 而由 (\SGARefLink{sga7.sga7i.eix.n11-6-7}{11.6.7}), 此核的泛纤维为单位群, 故其核为单位群 (\SGARefLink{sga7.sga7i.eix.n3-1-3}{3.1.3}). 因而 (\SGARefLink{sga7.sga7i.eix.n11-6-8-1}{11.6.8.1}) 是平展单态射, 从而为开浸入 (EGA IV 17.9.1), 引理得证.

现置

\begin{equation}\tag{\SGARefTarget{sga7.sga7i.eix.n11-6-9}{11.6.9}}
{}_n\Phi=({}_nA)^f/({}_nA^o)^f\subset{}_n\Psi.
\end{equation}
故 ${}_n\Phi$ 是 ${}_n\Psi$ 的有限开子概形, 因为 $({}_nA)^f$ 是 ${}_nA$ 的有限开子群概形. 更确切地, 立即有

\begin{equation}\tag{\SGARefTarget{sga7.sga7i.eix.n11-6-10}{11.6.10}}
{}_n\Phi\xrightarrow{\ \sim\ }({}_n\Psi)^f.
\end{equation}
将 (\SGARefLink{sga7.sga7i.eix.n11-6-9}{11.6.9}) 与 (\SGARefLink{sga7.sga7i.eix.n11-6-8-1}{11.6.8.1}) 复合, 得到开浸入

\begin{equation}\tag{\SGARefTarget{sga7.sga7i.eix.n11-6-11}{11.6.11}}
{}_n\Phi\lhook\joinrel\longrightarrow\underline M_n.
\end{equation}
另一方面,

\begin{equation*}
{}_n\Phi\times_Ss={}_n\Psi\times_Ss={}_nA_o/{}_nA_o^o,
\end{equation*}
且由于 $A_o^o$ 为 $n$-可除群 (\SGARefLink{sga7.sga7i.eix.n2-2-1}{2.2.1}), 最后一项立即等同于 ${}_n\Phi_o$. 因而得到同构

\begin{equation}\tag{\SGARefTarget{sga7.sga7i.eix.n11-6-12}{11.6.12}}
{}_n\Phi\times_Ss\xrightarrow{\ \sim\ }{}_n\Phi_o,
\end{equation}
这表明, ${}_n\Phi$ 在唯一同构的意义下, 被刻画为特殊纤维为 ${}_n\Phi_o$ 的 $S$ 上有限平展概形. 因而要定义所需的同构 (\SGARefLink{sga7.sga7i.eix.n11-6-4}{11.6.4}), 只需利用开浸入 (\SGARefLink{sga7.sga7i.eix.n11-6-11}{11.6.11}) 并证明以下结果.

\medskip
\noindent
\underline{引理} \SGARefTarget{sga7.sga7i.eix.n11-6-13}{11.6.13}. \underline{典范同态 (\SGARefLink{sga7.sga7i.eix.n11-6-11}{11.6.11}) 的像等于下列同态的核:}

\begin{equation*}
\widetilde u_n:\underline M_n\longrightarrow\underline{\check M}'_n.
\end{equation*}

由定义立即可知, 同态 (\SGARefLink{sga7.sga7i.eix.n11-6-11}{11.6.11}) 在 $n$ 变化时与过渡同态相容. 因此, 引理 \SGARefLink{sga7.sga7i.eix.n11-6-13}{11.6.13} 将完成定理 \SGARefLink{sga7.sga7i.eix.n11-5}{11.5} 的证明.

% END INCLUDED SGA7I_English_current_buildable_source_20260823/components/169_expose_IX_section_11_6_opening_through_lemma_11_6_13.tex

% BEGIN INCLUDED SGA7I_English_current_buildable_source_20260823/components/170_expose_IX_lemma_11_6_13_proof_and_remarks_11_7.tex
% R3 explicit revisions IX-B-04, IX-B-05, IX-B-06, IX-B-07, IX-B-08, IX-B-09, IX-B-10, R3-IX-B-15; see qa/FIDELITY-R3-CHANGES.json.
% Expose IX, proof of Lemma 11.6.13 and Remarks 11.7.
% French transcription lines 3304--3391; authority physical pp. 472--474
% (scan indices 471--473; source folios 460--462). Physical p. 474 was
% checked through the exact 11.8 boundary. All pages were compared in four
% overlapping 1100-dpi bands. The two compact diagrams, their nodes, ten plus
% five arrows, equality/isomorphism marks, labels, bars, and punctuation are
% all decisive at that resolution; no higher-resolution escalation was needed.

为证明引理, 注意

\begin{equation*}
K_n=\ker\widetilde u_n
\end{equation*}
也是下列典范配对的左零化子:

\begin{equation*}
u_n:\underline M_n\times\underline M'_n
 \longrightarrow(\mathbb Z/n\mathbb Z)_S,
\end{equation*}
其显式定义由 \SGARefLink{sga7.sga7i.eix.n9-4}{9.4} 及 \SGARefLink{sga7.sga7i.eix.n9-6}{9.6} 给出, 使用了扩张

\begin{equation}\tag{\SGARefTarget{sga7.sga7i.eix.n11-6-14}{11.6.14}}
\resizebox{\textwidth}{!}{%
\begin{tikzcd}[ampersand replacement=\&, column sep=small, row sep=large]
0 \arrow[r]
 \& {\ensuremath{\check{\underline M}'_n(1)}} \arrow[r]
 \& F_n \arrow[r] \arrow[d, equal]
 \& R_n \arrow[r]
 \& 0
 \& {,}
 \& 0 \arrow[r]
 \& {\ensuremath{\check{\underline M}_n(1)}} \arrow[r]
 \& F'_n \arrow[r] \arrow[d, equal]
 \& R'_n \arrow[r]
 \& 0 \\
{}
 \& {}
 \& ({}_nA^o)^f
 \& {}
 \& {}
 \& {}
 \& {}
 \& {}
 \& ({}_nA'^o)^f
 \& {}
 \& {}
\end{tikzcd}}
\end{equation}
(其中 $R_n$ 与 $R'_n$ 互为 Cartier 对偶), 以及$E_{n\eta}$ 被 $F_{n\eta}$ 的混合扩张 $X_{n\eta}$, 即 (\SGARefLink{sga7.sga7i.eix.n11-6-5}{11.6.5}) 的泛纤维;\footnote{印本此处引用了不存在的公式 (11.7.5); 这里取泛纤维的滤过是 (\SGARefLink{sga7.sga7i.eix.n11-6-5}{11.6.5}).} 此处置

\begin{equation}\tag{\SGARefTarget{sga7.sga7i.eix.n11-6-15}{11.6.15}}
\begin{tikzcd}[column sep=large, row sep=large]
E_n=D(F'_n),
 & 0 \arrow[r]
 & R_n \arrow[r] \arrow[d, no head, "\wr"']
 & E_n \arrow[r]
 & \underline M_n \arrow[r]
 & 0
 & {.} \\
{}
 & {}
 & D(R'_n)
 & {}
 & {}
 & {}
 & {}
\end{tikzcd}
\end{equation}
由命题 \SGARefLink{sga7.sga7i.eix.n9-4-7}{9.4.7}\footnote{印本引用 \SGARefLink{sga7.sga7i.eix.n9-4-6}{9.4.6}, 但所需的扩张判别是命题 \SGARefLink{sga7.sga7i.eix.n9-4-7}{9.4.7}.} 立即可知, 在 $\underline M_n$ 的所有有限平展子群 $K'_n$ 中, 若要求 $K'_n$ 在混合扩张 $X_{n\eta}$ 中的逆像 $\overline X_{n\eta}$, 作为 $\overline E_{n\eta}$ 被 $F_{n\eta}$ 的混合扩张, 可延拓为 $\overline E_n$ 被 $F_n$ 的混合扩张, 则 $K_n$ 是其中最大的. 这里 $\overline E_n$ 是显然的、以 $R_n$ 为核的 $K'_n$ 的子扩张.\footnote{印本在这一一般判别中写成“$K_n$ 的”, 但此处变化的子群是 $K'_n$; 直到下文才取特例 $K'_n=K_n$.}

(\SGARefLink{sga7.sga7i.eix.n11-6-11}{11.6.11}) 中 $\underline M_n$ 的子群 ${}_n\Phi$ 显然满足此条件, 因为有 $S$ 上相应的混合扩张 $({}_nA)^f$. 故 ${}_n\Phi\subset K_n$. 为证明反向包含, 考虑取 $K'_n=K_n$ 后按上述方式定义的、延拓 $\overline X_{n\eta}$ 的 $\overline E_n$ 被 $F_n$ 的混合扩张 $\overline X_n$. 因而它是扩张

\begin{equation*}
0\longrightarrow F_n\longrightarrow\overline X_n
 \longrightarrow K_n\longrightarrow0,
\end{equation*}
且有典范同态

\begin{equation}\tag{*}
F_n=({}_nA^o)^f\lhook\joinrel\longrightarrow A,
\end{equation}
其在泛纤维上由下列典范同态诱导:

\begin{equation}\tag{**}
\overline X_{n\eta}\lhook\joinrel\longrightarrow X_{n\eta}
 ={}_nA_K\lhook\joinrel\longrightarrow A_K.
\end{equation}
由 (\SGARefLink{sga7.sga7i.eix.n5-9-2}{5.9.2}) 立即可知, 同态 $(**)$ 唯一延拓为同态

\begin{equation}\tag{***}
\overline X_n\longrightarrow A
\end{equation}
诱导 $(*)$. 由于后一个同态将泛纤维映入 $A$ 的闭子概形 ${}_nA$, 它经过 $\overline X_n\to{}_nA$ 分解; 又由于 $\overline X_n$ 有限, 它经过 $({}_nA)^f$ 分解. 考察诱导的泛纤维上的同态, 得出

\begin{equation*}
\overline X_{n\eta}\subset({}_nA)^f_\eta,
\quad\text{即}\quad
K_{n\eta}\subset{}_n\Phi_\eta,
\quad\text{即}\quad
K_n\subset{}_n\Phi.
\end{equation*}
引理得证.

\medskip
\noindent
\underline{注} \SGARefTarget{sga7.sga7i.eix.n11-7}{11.7}. \SGARefLink{sga7.sga7i.eix.n11-6}{11.6} 中所作的构造\footnote{印本在 \SGARefLink{sga7.sga7i.eix.n11-7}{11.7} 开头写成“\SGARefLink{sga7.sga7i.eix.n11-7}{11.7} 中所作的构造”; 所指是刚在 \SGARefLink{sga7.sga7i.eix.n11-6}{11.6} 完成的构造.} 表明, 由 $S$ 上的扩张 (\SGARefLink{sga7.sga7i.eix.n11-6-14}{11.6.14}) 以及$E_n$ 被 $F_n$ 的混合扩张 $X_{n\eta}={}_nA_K$, 可以在唯一同构的意义下重构混合扩张 (\SGARefLink{sga7.sga7i.eix.n11-6-5}{11.6.5}), 从而也可重构 ${}_nA$ 的子群

\begin{equation*}
({}_nA)^f\supset({}_nA^o)^f
\end{equation*}
事实上, 由 \SGARefLink{sga7.sga7i.eix.n9-4}{9.4}, 这些资料决定配对

\begin{equation*}
u_n:\underline M_n\times\underline M'_n
 \longrightarrow(\mathbb Z/n\mathbb Z)_S,
\end{equation*}
而上述构造将 $({}_nA)^f$ 给出为 $\overline E_n$ 被 $F_n$ 的混合扩张 $\overline X_n$, 延拓了 $X_{n\eta}$ 诱导的混合扩张 $\overline X_{n\eta}$. 这里 $\overline E_n$ 是 $\ker\widetilde u_n$ 在 $E_n\to\underline M_n$ 下的逆像. 由 (9.4.8), 这一混合扩张 $\overline X_n$ 因而在\underline{唯一同构的意义下}被确定. 随后将 $\overline X_n$ 与 $X_{n\eta}$ 沿其公共开子概形 $\overline X_{n\eta}$ 粘合, 即恢复 ${}_nA$.

因此, 若已知结构资料 (\SGARefLink{sga7.sga7i.eix.n9-6-1}{9.6.1}) 和$E_\eta$ 被 $F_\eta$ 的混合扩张

\begin{equation*}
X_\eta=T_\ell(A_K)
\end{equation*}
——等价地, 对每个 $n=\ell^{r+1}$, $r\geq0$, 均已知扩张 (\SGARefLink{sga7.sga7i.eix.n11-6-14}{11.6.14}) 和 $X_{n\eta}$——则可重构混合扩张 (\SGARefLink{sga7.sga7i.eix.n11-6-5}{11.6.5}) 所成的归纳系统 (也可视为投射系统).

% END INCLUDED SGA7I_English_current_buildable_source_20260823/components/170_expose_IX_lemma_11_6_13_proof_and_remarks_11_7.tex

% BEGIN INCLUDED SGA7I_English_current_buildable_source_20260823/components/171_expose_IX_sections_11_8_through_11_11.tex
% Expose IX, sections 11.8--11.11.
% French transcription lines 3393--3446; authority physical pp. 474--476
% (scan indices 473--475; source folios 462--464). Physical p. 477 was
% checked through the exact section-12 boundary. All pages were compared in
% four overlapping 1100-dpi bands; every formula, accent, subscript, arrow,
% underline, and punctuation mark was decisive without higher escalation.
% The page images restore two transcription slips: the quotient sign in
% (11.8.1), and Phi rather than Psi on the left of (11.8.3).

\SGARefTarget{sga7.sga7i.eix.n11-8}{11.8}. 与 \SGARefLink{sga7.sga7i.eix.n8-3}{8.3} 一样, 将 $K$ 的可分闭包 $\overline K$ 视为其有限子扩张 $K'$ 的归纳极限. 对每个这样的 $K'$, 以 $A[K']$ 表示

\begin{equation*}
A_{K'}=A_K\otimes_K K'.
\end{equation*}
的 Néron 模型. 它给出连通分支的有限平展群概形

\begin{equation}\tag{\SGARefTarget{sga7.sga7i.eix.n11-8-1}{11.8.1}}
\Phi_o[K']=A[K']_o/A[K']_o^o
\end{equation}
它定义在 $S$ 在 $K'$ 中的正规化 $S[K']$ 的剩余域上. 为简便起见, 假设 $S$ 的剩余域可分闭, 因而各 $S[K']$ 的剩余域亦然. 此时群 (\SGARefLink{sga7.sga7i.eix.n11-8-1}{11.8.1}) 可等同于通常的有限群 $\Phi[K']$, 且显然当 $K'$ 变化时它们构成归纳系统, 下面就研究这一系统. 特别关心群

\begin{equation}\tag{\SGARefTarget{sga7.sga7i.eix.n11-8-2}{11.8.2}}
\Psi=\varinjlim_{K'}\Phi[K'].
\end{equation}
由 \SGARefLink{sga7.sga7i.eix.n3-6}{3.6} 知道, 当 $K'$ 足够大, 例如 $K'\supset K'_o$ 时, $A[K']$ 有半稳定约化. 因而当 $K'\supset K'_o$ 时, 过渡态射为单射 (\SGARefLink{sga7.sga7i.eix.n3-3}{3.3}). 此外, 本节沿用 \SGARefLink{sga7.sga7i.eix.n5-5}{5.5} 的记号, 当 $K'\supset K'_o$ 时有典范同构

\begin{equation}\tag{\SGARefTarget{sga7.sga7i.eix.n11-8-3}{11.8.3}}
\Phi[K'](\ell)\simeq
\ker\bigl(\widetilde u[K']_\ell\otimes D_\ell:
 M_\ell\otimes D_\ell\longrightarrow
 \check M'_\ell\otimes D_\ell\bigr),
\end{equation}
它来自将 (\SGARefLink{sga7.sga7i.eix.n11-6-2}{11.6.2}) 应用于 $A[K']$. 也可立即验证, 相应的浸入

\begin{equation*}
\Phi[K'](\ell)\lhook\joinrel\longrightarrow M_\ell\otimes D_\ell
\end{equation*}
当 $K'$ 变化时与过渡态射相容. 因而取极限给出 $\Psi$ 的 $\ell$-准素分量出发的典范单态射

\begin{equation}\tag{\SGARefTarget{sga7.sga7i.eix.n11-8-4}{11.8.4}}
\Psi(\ell)\mathrel{\mathop{=}^{\mathrm{def}}}
 \varinjlim_{K'}\Phi[K'](\ell)
 \lhook\joinrel\longrightarrow M_\ell\otimes D_\ell
\end{equation}
对所有 $\ell$ 求和, 得

\begin{equation}\tag{\SGARefTarget{sga7.sga7i.eix.n11-8-5}{11.8.5}}
\Psi\mathrel{\mathop{=}^{\mathrm{def}}}
 \varinjlim_{K'}\Phi[K']
 \longrightarrow M\otimes\mathbb Q/\mathbb Z.
\end{equation}

\medskip
\noindent
\underline{定理} \SGARefTarget{sga7.sga7i.eix.n11-9}{11.9}. \underline{典范同态 (\SGARefLink{sga7.sga7i.eix.n11-8-5}{11.8.5}) 为同构}.

已经知道它是单态射. 为证明它是同构, 必须使用 $\Phi[K'](\ell)$ 的像的刻画 (\SGARefLink{sga7.sga7i.eix.n11-8-3}{11.8.3}), 以及 $\widetilde u[K']$ 在 $K'$ 变化时的行为. 这种行为由 (\SGARefLink{sga7.sga7i.eix.n10-3-4}{10.3.4}) 给出.

% SGA-SRC-001281 — internal reference omitted in source; target not recoverable
可归结为 $S$ 严格局部的情形. 此时, 为证明 \SGARefLink{sga7.sga7i.eix.n11-9}{11.9}, 只需将下列熟知的断言应用于 $S[K'_o]$: 若 $S$ 为严格局部性状, 分式域为 $K$, 则对每个整数 $n\geq1$, 都存在 $K$ 的可分扩张 $K'$, 其次数可被 $n$ 整除.
\footnote{印本附有“参见 I”, 后面留空了内部引用. 原文没有给出目标, 此处也不另行编造.}

\medskip
\noindent
\underline{推论} \SGARefTarget{sga7.sga7i.eix.n11-10}{11.10}. \underline{在 \SGARefLink{sga7.sga7i.eix.n8-3}{8.3} 的假设 ($S$ 完备) 下并沿用其记号, 群概形 $A_K^{\flat}$ 的连通分支所成的 $K$ 上 ind 有限群概形, 即 $A_K^{\flat}/{A_K^{\flat}}^o$, 典范同构于} $\underline M_K\otimes\mathbb Q/\mathbb Z$, \underline{其中 $\underline M_K$ 为 (\SGARefLink{sga7.sga7i.eix.n8-1-8}{8.1.8}) 所定义的 $K$ 上的扭常群}.

由 $\underline M_K$ 的定义, 这立即来自 (8.2.4) 和 \SGARefLink{sga7.sga7i.eix.n11-9}{11.9}.

\medskip
\noindent
\underline{注} \SGARefTarget{sga7.sga7i.eix.n11-11}{11.11}. a) 若限于 $\ell\neq p$ 时的 $\ell$-准素分量, 则定理 \SGARefLink{sga7.sga7i.eix.n11-9}{11.9} 立即由 \SGARefLink{sga7.sga7i.eix.n11-2}{11.2} 得到, 无需 \SGARefLink{sga7.sga7i.eix.n11-6}{11.6} 中更精细的考察 (经由 \SGARefLink{sga7.sga7i.eix.n9-4}{9.4}). 同样, 若假设 $K$ 为\underline{零特征}, 则 \SGARefLink{sga7.sga7i.eix.n11-9}{11.9} 的证明无需使用 (\SGARefLink{sga7.sga7i.eix.n11-8-3}{11.8.3}): 显然 ${}_n\Phi[K'](\ell)$ 在

\begin{equation*}
M_n={}_n(M_\ell\otimes D_\ell)
\end{equation*}
中的像包含 ${}_nA$ 的 $K'$-值点群 ${}_nA(K')$ 的像. 当 $K'$ 足够大时, 此群等于 ${}_nA(\overline K)$, 因而满射到 $M_n$. 另一方面, 若 $\operatorname{char}K=p$, 看来不存在“平凡”的证明, 即不存在绕开 \SGARefLink{sga7.sga7i.eix.n9-4}{9.4} 的、对 $p$-准素分量的 \SGARefLink{sga7.sga7i.eix.n11-10}{11.10} 的证明.

b) 同构

\begin{equation*}
\Psi\simeq(\mathbb Q/\mathbb Z)^r,
\end{equation*}
的存在性, 连同 $r$ 的正确取值, 是 J.-P. Serre 于 1964 年提出的猜想.

% END INCLUDED SGA7I_English_current_buildable_source_20260823/components/171_expose_IX_sections_11_8_through_11_11.tex

% BEGIN INCLUDED SGA7I_English_current_buildable_source_20260823/components/172_expose_IX_section_12_1_through_12_2.tex
% R3 explicit revisions IX-B-11; see qa/FIDELITY-R3-CHANGES.json.
\medskip
\noindent
12. \underline{Jacobian 的 Néron 模型与 Picard--Lefschetz 公式}

先回顾 M. Raynaud 的下列一般结果 [22, 定理~5 和 7; 22 bis].

\medskip
\noindent
\underline{定理} \SGARefTarget{sga7.sga7i.eix.n12-1}{12.1} (M. Raynaud). \underline{设 $S$ 为严格局部性状, $f:X\to S$ 为 $S$ 上的紧合平坦概形, 其纤维为一维, 且满足}

\begin{equation}\tag{\SGARefTarget{sga7.sga7i.eix.n12-1-1}{12.1.1}}
 f_*\underline{\mathcal O}_X\simeq\underline{\mathcal O}_S,
\end{equation}

\noindent
\underline{并且 $X$ 的局部环均为析因环. 对 $X_o$ 的每个极大点 $x_i$, 置}

\begin{equation}\tag{\SGARefTarget{sga7.sga7i.eix.n12-1-2}{12.1.2}}
 d_i=\operatorname{length}\underline{\mathcal O}_{X_o,x_i},
 \qquad d_i^t=\mu_i d_i,
\end{equation}

\noindent
\underline{其中 $\mu_i$ 是 $k(x_i)$ 在 $k$ 上的根重数} (EGA IV, 4.7.4) \underline{(因而为 $k$ 的特征指数的幂). 设 $d$ (相应地 $d^t$) 为整数 $d_i$ (相应地 $d_i^t$) 的最大公因数; 故若 $k$ 完美, 则 $d=d^t$. 置}

\begin{equation}\tag{\SGARefTarget{sga7.sga7i.eix.n12-1-3}{12.1.3}}
 P=\underline{\operatorname{Pic}}_{X/S}:
   (\operatorname{Sch})_{/S}^{,o}\longrightarrow(\operatorname{Ab}).
\end{equation}

\noindent
a) \underline{考虑下列条件}:

\smallskip
\noindent
1) $d^t=1$.

\smallskip
\noindent
2) $X$ \underline{有一个次数为 $1$ 的零维闭链}.

\smallskip
\noindent
3) \underline{$X\times P^o$ 上存在可逆层 $L$, 使相应的同态}
$P^o_\eta\to\underline{\operatorname{Pic}}_{X_\eta/\eta}=P_\eta$
\underline{为典范包含}. (由 J.-P. Murre 的结果已知, $P_\eta$ 可由 $\eta$ 上的局部有限型群概形表示; 例如参见 SGA 6, XII, 1.5.)

\smallskip
\noindent
4) $X$ \underline{在 $S$ 上上同调平坦} (EGA III, 7.8.1) (等价地, 在零维上同调平坦; 又等价于 $k\xrightarrow{\sim}H^o(X_o,\underline{\mathcal O}_{X_o})$), \underline{且 $d=1$}.

\smallskip
\noindent
5) $P$ \underline{可表} (等价地, 可由 $S$ 上的光滑群概形表示).

\smallskip
\noindent
6) $P^o$ \underline{可表} (等价地, 可由 $S$ 上的光滑群概形表示). ($P^o$ 的定义见 SGA 3, $\mathrm{VI}_B$, \SGARefLink{sga7.sga7i.eix.n3-1}{3.1}.)

\smallskip
\noindent
7) $P^o$ \underline{在 $S$ 上分离}, 即满足 EGA II, 7.2.3 的分离性赋值判别 (等价地, $P^o$ 在 $S$ 上的每个满足 $s(\eta)=0$ 的截面 $s$ 均为零).

\smallskip
\noindent
8) $d=1$.

\smallskip
\noindent
\underline{则有如下蕴含关系}

\begin{equation}\tag{\SGARefTarget{sga7.sga7i.eix.n12-1-4}{12.1.4}}
 1)\Longleftrightarrow2)\Longrightarrow3)\Longrightarrow4)
 \Longleftrightarrow5)\Longleftrightarrow6)\Longleftrightarrow7)
 \Longrightarrow8).
\end{equation}

\noindent
\underline{特别地, 若 $d=d^t$, 例如 $k$ 完美时, 条件 1) 至 8) 等价}.

\medskip
\noindent
b) \underline{假设 $d^t=1$, 因而 $P$ 可由 $S$ 上的光滑群概形表示. 则 $P$ 的单位截面的概形闭包是一个在 $S$ 上平展的子群概形 $E$. 若 $I$ 为 $X_o$ 的既约不可约分支 $X_o^i$ 所成的集合, 则有同构}

\begin{equation}\tag{\SGARefTarget{sga7.sga7i.eix.n12-1-5}{12.1.5}}
 \mathbb Z^I\simeq\Gamma(S,E)
\end{equation}

\noindent
\underline{它将每个 $i\in I$ 对应于 $P$ 在 $S$ 上的截面 $S_i$ (在 $\eta$ 上为零), 该截面由 $X$ 上的可逆层 $\underline{\mathcal O}_X(X_o^i)$ 定义}\footnote{印本此处写成“在 $S$ 上”, 后面又将同一除子写作 $X_i^o$. 这个层是 $X$ 上与分支 $X_o^i$ 相伴的可逆层; 英文版始终使用一致的记号.}. (由于按假设 $X$ 的局部环为析因环, 闭链 $X_o^i$ 确为除子; 因此 $\underline{\mathcal O}_X(X_o^i)$ 确实具有上述含义, 即与这一除子相伴的可逆层; 见 EGA IV, 21.2.8.1.) \underline{有}

\begin{equation}\tag{\SGARefTarget{sga7.sga7i.eix.n12-1-6}{12.1.6}}
 P^o\cap E=\underline{P\text{ 的零截面}}.
\end{equation}

\noindent
c) \underline{在 b) 的条件下, 商}

\begin{equation}\tag{\SGARefTarget{sga7.sga7i.eix.n12-1-7}{12.1.7}}
 Q=P/E
\end{equation}

\noindent
\underline{可由 $S$ 上的光滑分离群概形表示, 且这个群概形是 Néron 的, 即它相对于其特殊纤维满足泛性质 (\SGARefLink{sga7.sga7i.eix.n1-1-2}{1.1.2})}.

\medskip
\noindent
d) \underline{在 c) 的条件下, 进一步假设 $X$ 的几何泛纤维整. 则存在典范正合列}

\begin{equation}\tag{\SGARefTarget{sga7.sga7i.eix.n12-1-8}{12.1.8}}
 0\longrightarrow A\longrightarrow Q
   \xrightarrow{\ \deg\ }\mathbb Z_S\longrightarrow0,
\end{equation}

\noindent
\underline{其中 $\deg$ 由“次数”同态 $P\to\mathbb Z_S$ 诱导, 而 $A$ 是 $S$ 上的有限型光滑分离群概形, 同样也是 Néron 的, 即它为下列群概形的 Néron 模型:}

\begin{equation}\tag{\SGARefTarget{sga7.sga7i.eix.n12-1-9}{12.1.9}}
 A_\eta=P_\eta^o=
 \underline{\operatorname{Pic}}^o_{X_\eta/\eta}.
\end{equation}

\noindent
\underline{此外, 复合同态 $P^o\to P\to Q$ 诱导同构}

\begin{equation}\tag{\SGARefTarget{sga7.sga7i.eix.n12-1-10}{12.1.10}}
 P^o\xrightarrow{\ \sim\ }A^o.
\end{equation}

\medskip
\noindent
\SGARefTarget{sga7.sga7i.eix.n12-1-11}{12.1.11}. 我们关心的情形是 $X_\eta$ 在 $\eta$ 上光滑的情形; 由 (\SGARefLink{sga7.sga7i.eix.n12-1-1}{12.1.1}), 此时 $X$ 的几何泛纤维光滑且连通 (从而整), 并且 $X_\eta$ 有一个次数为一的零维闭链. 于是 \SGARefLink{sga7.sga7i.eix.n12-1}{12.1} 所考虑的全部条件 (特别是 a) 中的条件 1) 至 8)) 均成立. 此外, 此时 $A_\eta=P_\eta^o$ 为阿贝尔概形, 而 d) 给出了其 Néron 模型 $A$ 的十分显式的构造. 特别地, (\SGARefLink{sga7.sga7i.eix.n12-1-10}{12.1.10}) 给出同构

\begin{equation}\tag{\SGARefTarget{sga7.sga7i.eix.n12-1-12}{12.1.12}}
 A_o^{,o}\simeq\underline{\operatorname{Pic}}^o_{X_o/k},
\end{equation}

\noindent
特别地, 这表明 $A$ \underline{在 $S$ 上有半稳定约化, 当且仅当} $\underline{\operatorname{Pic}}^o_{X_o/k}$ \underline{的幺幂秩为零}. 当 $X$ 正则且 $X_o$ 在 $k$ 上可分时, 已知这一条件成立, 当且仅当几何特殊纤维 $\overline X_o$ 除普通二重点外没有其他奇点 (即该点的完备局部环与仿射 $2$ 维空间中两条坐标轴之并在原点处的完备局部环相同). 结合域 $k$ 上紧合可分曲线的 Picard 概形的“熟知”描述, 利用 EGA IV, 21.8.5 验证这一事实, 留给读者作为练习. 还应注意, 利用 \SGARefLink{sga7.sga7i.eix.n3-5}{3.5} 中半稳定约化的 Galois 判别, 蕴含关系“$\overline X_o$ 具有普通二重点” $\Longrightarrow$ “$A_\eta=\underline{\operatorname{Pic}}^o_{X_\eta/\eta}$ 有半稳定约化”立即由“几何单值定理” III\footnote{印本中“III”后面的内部定位留空; 这里不另行编造目标.} 得到. 不过, 这一论证并不能给出重要的同构 (\SGARefLink{sga7.sga7i.eix.n12-1-10}{12.1.10}) 及其推论 (\SGARefLink{sga7.sga7i.eix.n12-1-12}{12.1.12}).

\medskip
\noindent
\SGARefTarget{sga7.sga7i.eix.n12-2}{12.2}. \SGARefLink{sga7.sga7i.eix.n12-1}{12.1} 的证明见 [22 bis]; 实际上, 我们只需 $X_\eta$ 在 $\eta$ 上光滑的情形. 还应指出, [22 bis] 在比 \SGARefLink{sga7.sga7i.eix.n12-1}{12.1} c) 和 d) 广得多的条件下证明了 Picard 概形与 Néron 模型之间的关系. 由于这些陈述技术性较强, 本综述不予列入. 只需指出, b)、c)、d) 容易由 a) 得出, 而 a) 中仅有的细致之处在于蕴含关系 $3)\Longrightarrow4)\Longrightarrow5)$ 和 $7)\Longrightarrow4)$, 其中第一个最为困难. 不过, 若假设 $X_o$ 在 $k$ 上可分, 这一点就变得平凡 (因为 $S$ 严格局部, 这意味着 $X$ 在 $S$ 上有截面, 从而条件 1), 即 $d^t=1$, 成立). 建议读者在这一特别简单的情形下给出证明, 必要时可再假设 $X_\eta$ 光滑, 从而 $P_\eta^o$ 为阿贝尔概形. 此时可使用其 Néron 模型 $A$, 且 $A^o$ 表示 $\underline{\operatorname{Pic}}^o_{X/S}$. 我们仅在这一特殊情形下使用 \SGARefLink{sga7.sga7i.eix.n12-1}{12.1}, 并进一步假设 $\overline X_o$ 除普通二重点外没有其他奇点, 即由 (\SGARefLink{sga7.sga7i.eix.n12-1-11}{12.1.11}), $A_\eta$ 有半稳定约化. 这是 12.5.4 所考虑的情形, 从现在起就限于此情形.

% END INCLUDED SGA7I_English_current_buildable_source_20260823/components/172_expose_IX_section_12_1_through_12_2.tex

% BEGIN INCLUDED SGA7I_English_current_buildable_source_20260823/components/173_expose_IX_section_12_3.tex
% R3 explicit revisions IX-B-12; see qa/FIDELITY-R3-CHANGES.json.
\medskip
\noindent
\SGARefTarget{sga7.sga7i.eix.n12-3}{12.3}. 设 $C$ 为域 $k$ 上的紧合可分曲线; 为简便起见, 假设它可分闭 (读者也可假设其代数闭). 若 $\widetilde C$ 为 $C$ 的正规化, 已知 $\underline{\operatorname{Pic}}^o_{\widetilde C/k}$ 的“约化秩”为零 (即其极大环面为零). 因而 $\underline{\operatorname{Pic}}^o_{C/k}$ 的极大环面包含于

\begin{equation*}
 N=\ker\bigl(\underline{\operatorname{Pic}}_{C/k}
       \longrightarrow\underline{\operatorname{Pic}}_{\widetilde C/k}\bigr),
\end{equation*}

\noindent
因而典范同构于后一个群的极大环面. 利用 EGA IV, 21.8.5 容易显式给出 $N$ 的结构: $N$ 典范地为以连通幺幂群 $U$ 为核的某个环面 $T$ 的扩张, 我们将描述这个环面 (该连通幺幂群本身在 $k$ 上有一个合成列, 其商同构于加法群 $\mathbb G_a$). 此外, 若 $T'$ 为 $N$ 的极大环面, 则 $T'\cap U$ 既幺幂又为乘法型, 因而为零 (SGA 3, XVII, 2.4), 而 $T'$ 在 $T$ 中的像为整个 $T$ (SGA 3, XII, 7.1 e)). 故 $T'\to T$ 为同构. 因此 $T$ 典范同构于 $N$ 的极大环面, 从而同构于 $\underline{\operatorname{Pic}}^o_{C/k}$ 的极大环面.
\footnote{印本在这段论证末尾写成 $\underline{\operatorname{Pic}}^o_{X/k}$. 整个 \SGARefLink{sga7.sga7i.eix.n12-3}{12.3} 中的曲线均为 $C$, 前后的相应记号均为 $\underline{\operatorname{Pic}}^o_{C/k}$; 英文版恢复为 $C$.}

为显式给出 $N$ 的结构, 方便起见假设 $k$ 可分闭. 此时 $C$ 的不可约分支 $C_i$ 几何不可约, 且 $C$ 的非正规点在 $k$ 上为根的. 设 $J$ 为 $\widetilde C$ 的不可约分支集合; 对每个 $x\in C$, 设 $J(x)$ 为 $C$ 在 $x$ 处的“分枝”集合 (即 $\widetilde C$ 中位于 $x$ 上方的点). 最后置

\begin{equation}\tag{\SGARefTarget{sga7.sga7i.eix.n12-3-1}{12.3.1}}
 R(x)=\mathbb Z^{J(x)}/\operatorname{Im}\mathbb Z,
\end{equation}

\noindent
其中 $\mathbb Z\to\mathbb Z^{J(x)}$ 为对角映射, 以及

\begin{equation}\tag{\SGARefTarget{sga7.sga7i.eix.n12-3-2}{12.3.2}}
 R=\left(\bigoplus_x R(x)\right)/\operatorname{Im}\mathbb Z^J.
\end{equation}

\noindent
显然只需对 $C$ 中同时属于 $C$ 的至少两条分枝的点求和, 这些点构成有限集. 因而 $R$ 为有限生成 $\mathbb Z$-模. 此时有典范同构

\begin{equation}\tag{\SGARefTarget{sga7.sga7i.eix.n12-3-3}{12.3.3}}
 T\ \bigl(\simeq\text{ 为下列群的极大环面: }
       \underline{\operatorname{Pic}}^o_{C/k}\bigr)
 \simeq R\otimes\mathbb G_{m,k}.
\end{equation}

\noindent
换言之, $T$ 的特征标群 $M$ 为

\begin{equation}\tag{\SGARefTarget{sga7.sga7i.eix.n12-3-4}{12.3.4}}
 M\mathrel{\mathop{=}^{\mathrm{def}}}D(T)=\check R
 \simeq\ker\left(\bigoplus_x\check R(x)\longrightarrow\mathbb Z^J\right),
\end{equation}

\noindent
其中 $\check R(x)$ 也可解释为 $J(x)$ 上取值之和为零的整数值函数所成的集合.
\footnote{印本将这两个公式编号为 (12.2.3) 和 (12.2.4). 它们属于 \SGARefLink{sga7.sga7i.eix.n12-3}{12.3}, 下文也将其引为 (\SGARefLink{sga7.sga7i.eix.n12-3-3}{12.3.3}) 及 (\SGARefLink{sga7.sga7i.eix.n12-3-4}{12.3.4}); 此处采用这些一致的编号.}

\medskip
\noindent
\SGARefTarget{sga7.sga7i.eix.n12-3-5}{12.3.5}. 我们只需考虑 $C$ 的任意一点 $x$ 处至多有两条分枝的情形 (特别地, 普通二重点就是如此). 若 $\operatorname{card}J(x)=2$, 则 $R(x)$ 同构于 $\mathbb Z$, 但同构依赖于在 $x$ 处的两条分枝 $(C',C'')$ 的一个次序, 或等价地依赖于选定一条分枝 $C'$: 将 $C'$ 对应于 $R(x)$ 中的元素, 即 $0_{C'}+1_{C''}$ 的典范像. 此时得到同构

\begin{equation}\tag{\SGARefTarget{sga7.sga7i.eix.n12-3-6}{12.3.6}}
 R=\mathbb Z^I/\operatorname{Im}\mathbb Z^J,
 \qquad
 T\simeq\mathbb G_m^I/\operatorname{Im}\mathbb G_m^J,
 \qquad
 M=\ker\bigl(\mathbb Z^I\longrightarrow\mathbb Z^J\bigr),
\end{equation}

\noindent
其中 $I$ 为属于 $C$ 的两条分枝的点 $x\in C$ 所成的集合, 并已按刚才说明的方式选定符号.

\medskip
\noindent
\underline{注} \SGARefTarget{sga7.sga7i.eix.n12-3-7}{12.3.7}. 有时可方便地用如下方式描述 $\underline{\operatorname{Pic}}^o_{C/k}$ 的极大环面 $T$ 的特征标群 $M$. 考虑“无向一维图” $\Gamma(C)$ (N. Bourbaki, \emph{Groupes et alg\`ebres de Lie}, 第 IV 章, 附录), 其顶点为 $\widetilde C$ 的不可约分支. 两个顶点 $\widetilde C_i$ 与 $\widetilde C_j$ 之间的边集合, 与 $C$ 的如下成对分枝一一对应: 两条分枝\underline{互不相同}, 分别包含在 $\widetilde C_i$ 和 $\widetilde C_j$ 中, 且通过 $C$ 的同一点. 当 $i\neq j$ 时, 这个集合为 $\widetilde C_i\times_C\widetilde C_j$; 当 $i=j$ 时, 则为

\begin{equation*}
 \bigl(\widetilde C_i\times_C\widetilde C_i-
       \Delta(\widetilde C_i)\bigr)/(\text{交换对称}),
\end{equation*}

\noindent
其中 $\Delta:\widetilde C\to\widetilde C\times_C\widetilde C$ 为对角映射. (因而若 $C_i$ 几何单枝, 则该集合为空.) 此时有典范同构

\begin{equation}\tag{\SGARefTarget{sga7.sga7i.eix.n12-3-7.occ-2}{12.3.7}}
 M=D(T)\xleftarrow{\ \sim\ }H_1(\Gamma(C),\mathbb Z),
 \qquad
 \check M=R\simeq H^1(\Gamma(C),\mathbb Z).
\end{equation}

我们并不需要 $M$ 的这一组合解释, 仅定义典范同态

\begin{equation}\tag{\SGARefTarget{sga7.sga7i.eix.n12-3-8}{12.3.8}}
 H_1(\Gamma(C),\mathbb Z)\longrightarrow M,
\end{equation}

\noindent
将它为同构的证明留给读者. 设 $I$ 为 $C$ 中属于 $C$ 的至少两条分枝的点所成的集合, $J$ 如上为 $\widetilde C$ (或等价地 $C$) 的不可约分支集合, $K$ 为 $C$ 在 $I$ 的各点处的“分枝”集合, 即 $I$ 在 $\widetilde C$ 中的逆像. 有典范映射

\begin{equation}\tag{\SGARefTarget{sga7.sga7i.eix.n12-3-9}{12.3.9}}
 h:K\longrightarrow I\times J,
 \qquad h=(f,g),
\end{equation}

\noindent
而 $M$ 的描述 (\SGARefLink{sga7.sga7i.eix.n12-3-4}{12.3.4}) 可用 $h$ 表示为

\begin{equation}\tag{\SGARefTarget{sga7.sga7i.eix.n12-3-10}{12.3.10}}
 M\simeq\ker\left(
  \mathbb Z^{(h)}:\mathbb Z^{(K)}\longrightarrow
  \mathbb Z^{(I\times J)}\right).
\end{equation}

事实上, $R=\check M$ 的表达式 (\SGARefLink{sga7.sga7i.eix.n12-3-3}{12.3.3}) 显然可解释为 $h$ 诱导的同态 $\mathbb Z^{I\times J}\to\mathbb Z^K$ 的余核, 此同态为 $\mathbb Z^{(h)}$ 的转置. 图 $\Gamma(C)$ 同样可用 $h$ 描述为

\begin{equation*}
 \Gamma_0=J,
 \qquad
 \Gamma_1=(K\times_I K-\Delta(K))/(\text{交换对称}).
\end{equation*}

在 $\Gamma(C)$ 的每条边上选定一个定向, 以 $\Gamma_1$ 表示这些有向边 $(x,C',C'')$ 的集合. 则

\begin{equation}\tag{\SGARefTarget{sga7.sga7i.eix.n12-3-11}{12.3.11}}
 H_1(\Gamma(C),\mathbb Z)=
 \ker\left(d:\mathbb Z^{(\Gamma_1)}\longrightarrow\mathbb Z^{(J)}\right),
\end{equation}

\noindent
其中微分为

\begin{equation}\tag{\SGARefTarget{sga7.sga7i.eix.n12-3-12}{12.3.12}}
 d(x,C',C'')=g(C'')-g(C').
\end{equation}

同态 (\SGARefLink{sga7.sga7i.eix.n12-3-8}{12.3.8}) 由下列链复形同态 $v$ 诱导:\footnote{印本定义 (\SGARefLink{sga7.sga7i.eix.n12-3-8}{12.3.8}) 的公式虽然被后文引用, 却没有编号; 而 (\SGARefLink{sga7.sga7i.eix.n12-3-13}{12.3.13}) 将边链群 $\mathbb Z^{(\Gamma_1)}$ 写成 $\mathbb Z^{(\Gamma)}$. 英文版补上被引用的编号及所需的下标. 工作转录稿在自环边描述中写成的 $\widetilde C_i\times_C\widetilde C_1$, 也依照扫描件恢复为 $\widetilde C_i\times_C\widetilde C_i$.}

\begin{equation}\tag{\SGARefTarget{sga7.sga7i.eix.n12-3-13}{12.3.13}}
 \begin{split}
 v=(v_0,v_1):C_\bullet(\Gamma(C))
   &=\left[\mathbb Z^{(\Gamma_1)}\longrightarrow\mathbb Z^{(J)}\right]\\
   &\longrightarrow
     \left[\mathbb Z^{(K)}\longrightarrow\mathbb Z^{(I\times J)}\right],
 \end{split}
\end{equation}

\noindent
它定义为

\begin{equation}\tag{\SGARefTarget{sga7.sga7i.eix.n12-3-14}{12.3.14}}
 v_1(x,C',C'')=(x,C'')-(x,C'),
 \qquad
 v_0(\widetilde C_j)=(0,\widetilde C_j).
\end{equation}

% END INCLUDED SGA7I_English_current_buildable_source_20260823/components/173_expose_IX_section_12_3.tex

% BEGIN INCLUDED SGA7I_English_current_buildable_source_20260823/components/174_expose_IX_section_12_4_and_theorem_12_5.tex
\medskip
\noindent
\SGARefTarget{sga7.sga7i.eix.n12-4}{12.4}. 现在考虑 \SGARefLink{sga7.sga7i.eix.n12-2}{12.2} 末尾的情形. 将 (\SGARefLink{sga7.sga7i.eix.n12-1-12}{12.1.12}) 与 (\SGARefLink{sga7.sga7i.eix.n12-3-6}{12.3.6}) 应用于 $C=X_o$, 可得 $A=\underline{\operatorname{Pic}}^o_{X_\eta/\eta}$ 的 Néron 模型 $A$ 的特殊纤维 $A_o$ 的极大环面 $T_o$ 的特征标群 $M$ 为

\begin{equation}\tag{\SGARefTarget{sga7.sga7i.eix.n12-4-1}{12.4.1}}
 M=\ker\left(\mathbb Z^I\xrightarrow{\ d\ }\mathbb Z^J\right),
\end{equation}

\noindent
其中 $I$ 为 $X_o$ 的奇点集合, $J$ 为其不可约分支集合, 同态 $\mathbb Z^I\to\mathbb Z^J$ 给定为

\begin{equation}\tag{\SGARefTarget{sga7.sga7i.eix.n12-4-2}{12.4.2}}
 d(x_i)=C'_i-C''_i.
\end{equation}

这里 $C'_i$ 和 $C''_i$ 为 $\widetilde C$ 中与经过二重点 $x_i$ 的两条分枝相对应的两个不可约分支. 算子 $d$ 依赖于这两条分枝的一个次序 $(C'_i,C''_i)$ 的选取, 同构 (\SGARefLink{sga7.sga7i.eix.n12-4-1}{12.4.1}) 显然依赖于这样的一组选取. 不过, 对 $i\in I$, 若以

\begin{equation}\tag{\SGARefTarget{sga7.sga7i.eix.n12-4-3}{12.4.3}}
 \varphi_{x_i}:M\longrightarrow\mathbb Z
\end{equation}

\noindent
表示复合

\begin{equation*}
 M\longrightarrow\mathbb Z^I
   \xrightarrow{\ \operatorname{pr}_i\ }\mathbb Z,
\end{equation*}

\noindent
其中第一个同态来自 (\SGARefLink{sga7.sga7i.eix.n12-4-1}{12.4.1}), 则回到表达式 (\SGARefLink{sga7.sga7i.eix.n12-3-4}{12.3.4}) 后立即可见, $\varphi_{x_i}$ 在相差符号的意义下良定义. 更确切地, 它只依赖于所选的经过 $x_i$ 的两条分枝 $C'_i,C''_i$ 的次序, 在次序反转时变号. 特别地, 可知双线性型

\begin{equation}\tag{\SGARefTarget{sga7.sga7i.eix.n12-4-4}{12.4.4}}
 \varphi_{x_i}\otimes\varphi_{x_i}
 \in \check M\otimes_{\mathbb Z}\check M
 \simeq\operatorname{Hom}(M\otimes M,\mathbb Z)
\end{equation}

\noindent
由奇点 $x_i$ 典范决定, 没有任何符号歧义. 因而可置

\begin{equation}\tag{\SGARefTarget{sga7.sga7i.eix.n12-4-5}{12.4.5}}
 u=\sum_i\varphi_{x_i}\otimes\varphi_{x_i}:
 M\otimes_{\mathbb Z}M\longrightarrow\mathbb Z.
\end{equation}

这给出良定义的配对, 与符号的任何选取无关. 它显然\underline{正定}, 因为它是 $\mathbb Z^I$ 上的标准二次型 $\sum_i X_i^2$ 在 $M\hookrightarrow\mathbb Z^I$ 上的限制.

以 $M$ 表示 $T_o$ 的特征标群这一记号, 先验而言与前几节的记号并不一致: 前面将这个群记为 $M'$, 而以 $M$ 表示与 $A_K$ 的对偶阿贝尔概形 $A'_K$ 相伴的类似的群. 但由于 $A_K$ 为 Jacobian, 它有典范极化 $\xi$ [5]\footnote{印本在典范极化后的引文为“[5]”, 后面的具体定位留空. 英文版保留引文, 并未编造缺失的定位.}, 该极化诱导典范同构

\begin{equation}\tag{\SGARefTarget{sga7.sga7i.eix.n12-4-6}{12.4.6}}
 \psi_\xi:A_K\xrightarrow{\ \sim\ }A'_K.
\end{equation}

\noindent
利用此同构将 $A_K$ 与 $A'_K$ 等同, 从而无需再区分 $M$ 与 $M'$. 特别地, 对每个素数 $\ell$, 单值配对 (\SGARefLink{sga7.sga7i.eix.n9-1-2}{9.1.2}) 可视为配对

\begin{equation}\tag{\SGARefTarget{sga7.sga7i.eix.n12-4-7}{12.4.7}}
 u_\ell:M_\ell\otimes_{\mathbb Z_\ell}M_\ell
 \longrightarrow\mathbb Z_\ell.
\end{equation}

明确上述约定后, 可陈述本节的主要结果:

\medskip
\noindent
\underline{定理} \SGARefTarget{sga7.sga7i.eix.n12-5}{12.5} (\underline{Picard--Lefschetz 公式}).
\underline{设 $S$ 为严格局部性状, $X$ 为 $S$ 上的射影平坦正则概形, 其纤维为一维, 几何泛纤维光滑且连通, 几何特殊纤维除普通二次奇点外没有其他奇点. 沿用 \SGARefLink{sga7.sga7i.eix.n12-4}{12.4} 引入的记号, 对每个素数 $\ell$, 单值配对 (\SGARefLink{sga7.sga7i.eix.n12-4-7}{12.4.7}) 等于将 (\SGARefLink{sga7.sga7i.eix.n12-4-5}{12.4.5}) 的标量从 $\mathbb Z$ 扩张到 $\mathbb Z_\ell$ 所得的配对.}

% END INCLUDED SGA7I_English_current_buildable_source_20260823/components/174_expose_IX_section_12_4_and_theorem_12_5.tex

% BEGIN INCLUDED SGA7I_English_current_buildable_source_20260823/components/175_expose_IX_sections_12_6_and_12_7.tex
\medskip
\noindent
\SGARefTarget{sga7.sga7i.eix.n12-6}{12.6}. 注意, 在所考虑的特殊情形下, 若在 \SGARefLink{sga7.sga7i.eix.n10-4}{10.4} b) 中以 $\xi$ 为那里考虑的典范极化, 则这一陈述证明了 \SGARefLink{sga7.sga7i.eix.n10-4}{10.4}. 如 \SGARefLink{sga7.sga7i.eix.n10-5}{10.5} 中所示, 这足以推出一般情形下的 \SGARefLink{sga7.sga7i.eix.n10-4}{10.4}.

\medskip
\noindent
\SGARefTarget{sga7.sga7i.eix.n12-7}{12.7}. \underline{\SGARefLink{sga7.sga7i.eix.n12-5}{12.5} 的证明. $\ell\ne p$ 的情形.} 证明主要依赖于 \SGARefLink{sga7.sga7i.eix.n12-5}{12.5} 所考虑情形下的\underline{上同调 Picard--Lefschetz 公式}, 该公式给出单值群 $I$ 在下列群上的作用:\footnote{工作转录稿在 (\SGARefLink{sga7.sga7i.eix.n12-7-1}{12.7.1}) 的 $T_\ell(A_{\bar\eta})$ 上添加了对偶上标. 直接以 5,000-dpi 查看印本公式, 未见这一上标, 这也与标准等同 $H^1(A_{\bar\eta},\mathbb Z_\ell)=
\operatorname{Hom}(T_\ell(A_{\bar\eta}),\mathbb Z_\ell)$ 一致.}

\begin{equation}\tag{\SGARefTarget{sga7.sga7i.eix.n12-7-1}{12.7.1}}
 H^1(X_{\bar\eta},\mathbb Z_\ell)
 \simeq H^1(A_{\bar\eta},\mathbb Z_\ell)
 \simeq\operatorname{Hom}
   \bigl(T_\ell(A_{\bar\eta}),\mathbb Z_\ell\bigr)
\end{equation}

\noindent
这里 $\ell\ne p$. 这一公式将在以后的报告 (第 XV 篇) 中证明, 而且是在广得多的情形下证明; 它实际上等同于对所论素数 $\ell$ 的断言 \SGARefLink{sga7.sga7i.eix.n12-5}{12.5}, 因为 \SGARefLink{sga7.sga7i.eix.n9-2}{9.2} 在 $\ell\ne p$ 时给出了 $u_\ell$ 的定义. 因而只需处理 $\ell=p$ 的情形. 我们将通过一种将其归结为零特征 (并且使 $X_o$ 仅有一个奇点) 的方法来处理, 下面解释这一方法.

\medskip
\noindent
\underline{注 \SGARefTarget{sga7.sga7i.eix.n12-7-2}{12.7.2}.} 还应注意, 一旦完成上述归约, 通常的论证便可进一步将底域归为复数域 $\mathbb C$, 然后使用经典形式的 Picard--Lefschetz 公式——只是看来还没有令人满意的文献可供引用. 本讨论班给出的 Picard--Lefschetz 公式的证明也将以类似方式进行, 即归结到可用真正超越方法处理的超越情形. 因而可以说, 无论如何, 这里给出的 \SGARefLink{sga7.sga7i.eix.n12-5}{12.5} (从而 \SGARefLink{sga7.sga7i.eix.n10-4}{10.4}) 的证明在本质上是超越的. 还应指出, \SGARefLink{sga7.sga7i.eix.n10-4}{10.4} 的另一种已知证明来自刚性解析理论 [24], 因而同样应视为本质上“超越的”. 找到这些结果的纯代数证明将很有意义 (尽管其可能性远非显然!).

% END INCLUDED SGA7I_English_current_buildable_source_20260823/components/175_expose_IX_sections_12_6_and_12_7.tex

% BEGIN INCLUDED SGA7I_English_current_buildable_source_20260823/components/176_expose_IX_section_12_8_opening_through_before_12_8_9.tex
% R3 explicit revisions IX-B-13; see qa/FIDELITY-R3-CHANGES.json.
\medskip
\noindent
\SGARefTarget{sga7.sga7i.eix.n12-8}{12.8}. \underline{\SGARefLink{sga7.sga7i.eix.n12-5}{12.5} 的证明: 归结为 $S$ 为零特征的情形.} 所用论证的原理与第 V 篇中“单值群在 $\pi_1$ 上的作用本质温和的定理”的证明相同.\footnote{印本中“V”后面的内部定位留空, 没有给出节目标.} 注意, 由 \SGARefLink{sga7.sga7i.eix.n10-3}{10.3}, 可假设 $S$ 完备. 设 $W$ 为剩余域为 $k$ 的 Cohen $p$-环, 选取在剩余域上诱导恒等的同态 $W\longrightarrow V$ (EGA $0_{\mathrm{IV}}$, 19.8.6). 设 $(x_\lambda)_{\lambda\in L}$ 为 $X_o$ 的奇点集合, 考虑 $X_o$ 在 $W$ 上的形式模概形 $\underline S$, 连同“泛曲线”
\[
 \underline f:\underline X\longrightarrow\underline S
\]
其特殊纤维为 $X_o$ (第 VI 篇).\footnote{印本中“VI”后面的内部定位同样留空.} 在该处已知, $\underline S$ 同构于幂级数环 $W[[T_1,\ldots,T_N]]$ 的形式谱, 且在 $\underline S$ 中可找到正规交叉除子

\begin{equation}\tag{\SGARefTarget{sga7.sga7i.eix.n12-8-1}{12.8.1}}
 D=\sum_{\lambda\in L}D_\lambda,
\end{equation}

\noindent
其不可约分支 $D_\lambda$ 以 $L$ 为指标, 每个 $D_\lambda$ 由形如 $T_{i_\lambda}=0$ 的方程描述. (也可用整数集 $[1,r]$ 为 $L$ 编号, 从而假设 $i_\lambda=\lambda$.) 分支 $D_\lambda$ 是 $\underline f$ 的非光滑轨迹中经过 $X_o$ 的奇点 $x_\lambda$ 的分支 $D'_\lambda$ 在 $\underline f:\underline X\longrightarrow\underline S$ 下的像; 该轨迹为
\[
 D'=\bigcup_{\lambda\in L}D'_\lambda
\]
按构造本身, $f:X\longrightarrow S$ 同构于将 $\underline f:\underline X\longrightarrow\underline S$ 沿适当的 $W$-态射 $g:S\longrightarrow\underline S$ 作回拉所得:

\begin{equation}\tag{\SGARefTarget{sga7.sga7i.eix.n12-8-2}{12.8.2}}
\begin{tikzcd}
 \underline X \arrow[d,"\underline f"'] &
 X \arrow[l] \arrow[d,"f"] \\
 \underline S & S \arrow[l,"g"']
\end{tikzcd}
\qquad .
\end{equation}

\noindent
置

\begin{equation}\tag{\SGARefTarget{sga7.sga7i.eix.n12-8-3}{12.8.3}}
 P=\underline{\operatorname{Pic}}^{o}_{\underline X/\underline S}.
\end{equation}

\noindent
由于 $\underline f$ 的特殊纤维几何既约且维数为 $1$, 由 EGA III, 7.8.1 可知 $\underline f$ 在零维上同调平坦,\footnote{工作转录稿此处写作 $\underline F$. 直接查看印本页面可见是 $\underline f$, 而且本句所指为泛曲线的特殊纤维, 也要求如此.} 且 $P$ 在 $\underline S$ 上形式光滑 (后者由 $H^2(X_o,\underline{\mathcal O}_{X_o})=0$ 得出, 因为 $X_o$ 的维数为 $1$). 为简便起见, 假设 $P$ \underline{可由 $\underline S$ 上的群概形表示}; 它必为纤维连通的光滑群概形. 这是 [22] 引述的 Mumford 的未发表结果的特例 (Raynaud 在那里所引的札记中, 当底为性状时大大推广了这一结果). 由于 Mumford 和 Raynaud 的结果均无已发表的证明, 读者也可改用 M. Artin [1] 的一般结果. 这些结果仅要求 $\underline f$ 紧合、平坦且在零维上同调平坦, 便给出 $P$ 可由 Artin 意义下的“代数空间”表示, 这已足以满足我们的目的. 置

\begin{equation}\tag{\SGARefTarget{sga7.sga7i.eix.n12-8-4}{12.8.4}}
 U=\underline S-\operatorname{supp}D.
\end{equation}

\noindent
故 $U$ 为 $\underline S$ 中使 $\underline X$ 在其上光滑的最大开子集, 从而
\[
 P|_U=\underline{\operatorname{Pic}}^{o}_{\underline X_U/U}
\]
为 $U$ 上的阿贝尔概形. 它沿 $g:\eta\longrightarrow U$ 回拉到 $\eta$ 上, 就是我们要研究的 $\eta$ 上的阿贝尔概形 $A_\eta$. 更确切地, 有 $S$ 上的典范同构 (\SGARefLink{sga7.sga7i.eix.n12-1-10}{12.1.10}):

\begin{equation}\tag{\SGARefTarget{sga7.sga7i.eix.n12-8-5}{12.8.5}}
 P\times_{\underline S}S\simeq A^{o}.
\end{equation}

\noindent
由于特殊纤维 $P_o=A_o^{o}$ 对每个 $\ell$ 均为 $\ell$-可除群, $P$ 本身亦然; 因而其所有纤维都是以环面为核的阿贝尔簇的扩张. 考虑 $\underline S$ 上的 pro-$\ell$ 群 $T_\ell(P)$, 其沿 $g$ 的回拉恰为 $T_\ell(A^{o})$. 应用 \SGARefLink{sga7.sga7i.eix.n7-3}{7.3} 的构造, 可用“不变部分”和“环面部分”定义这个 pro 群的滤过:

\begin{equation}\tag{\SGARefTarget{sga7.sga7i.eix.n12-8-6}{12.8.6}}
 T_\ell(P)\supset T_\ell(P)^f=T_\ell(\widehat P)
 \supset T_\ell(P)^t=T_\ell(T)\supset0.
\end{equation}

\noindent
其沿 $g$ 的回拉为第 2 节与第 5 节研究过的 $T_\ell(A^{o})$ 的相应滤过.

考虑将此滤过限制到 $\underline S$ 的开子集 $U$ 上. 由于 $P_U$ 是 $U$ 上光滑曲线概形的 Jacobian, 它有典范自对偶, 从而有相应配对

\begin{equation}\tag{\SGARefTarget{sga7.sga7i.eix.n12-8-7}{12.8.7}}
 \varphi:T_\ell(P_U)\otimes T_\ell(P_U)
 \longrightarrow T_\ell(\mathbb G_{mU}).
\end{equation}

\noindent
这是 $T_\ell(P_U)$ 在 Barsotti--Tate 群意义下的\underline{自对偶}.\footnote{印本此处为 $T_\ell(B_U)$, 但没有引入过 $B_U$, 配对及所有相邻的项均涉及 $P_U$. 英文版恢复为上下文所要求的 $P_U$.} 在 \SGARefLink{sga7.sga7i.eix.n7-4-8}{7.4.8} 中已知, 对这一配对, $T_\ell(P_U)^f$ \underline{与} $T_\ell(P_U)^t$ \underline{互为正交补}. 因而
\[
 E_U=T_\ell(P_U)/T_\ell(P_U)^t
\]
典范同构于 $T_\ell(P_U)^f=F_U$ 的对偶 Barsotti--Tate pro 群. 因而它典范地延拓为 $\underline S$ 上的 Barsotti--Tate 群 $E$, 即

\begin{equation}\tag{\SGARefTarget{sga7.sga7i.eix.n12-8-8}{12.8.8}}
 E=D(F),\qquad\text{其中}\qquad F=T_\ell(P)^f.
\end{equation}

\noindent
现在处于典型的混合扩张情形 (\SGARefLink{sga7.sga7i.eix.n9-3}{9.3}), 可以重复 \SGARefLink{sga7.sga7i.eix.n9-4}{9.4}--\SGARefLink{sga7.sga7i.eix.n9-6}{9.6} 的构造 (这些构造先前写在将底概形限于性状的陈旧而寒酸的框架内, 谨为此致歉!): $T_\ell(P_U)$ 是 $E_U$ 被 $F_U$ 的混合扩张, 其中 $E$ 和 $F$ 为 $\underline S$ 上的 Barsotti--Tate 群, 分别为如下互为对偶的扩张:

% END INCLUDED SGA7I_English_current_buildable_source_20260823/components/176_expose_IX_section_12_8_opening_through_before_12_8_9.tex

% BEGIN INCLUDED SGA7I_English_current_buildable_source_20260823/components/177_expose_IX_section_12_8_9_through_12_10.tex
\begin{equation}\tag{\SGARefTarget{sga7.sga7i.eix.n12-8-9}{12.8.9}}
\begin{tikzcd}[row sep=small,column sep=small]
0 \arrow[r] & P \arrow[rr] \arrow[d,dash,"\wr"]
 & & F \arrow[r] \arrow[d,dash,"\wr"] & R \arrow[r] & 0 & \\
& T_\ell(P)^t\simeq
  \underline{\operatorname{Hom}}
    \bigl(M,T_\ell(\mathbb G_{m\underline S})\bigr)
& & T_\ell(P)^f & & & \\
& & & & & & \\
0 \arrow[r] & R \arrow[rr] \arrow[d,dash,"\wr"]
& & E \arrow[r] & Q \arrow[r] \arrow[d,dash,"\wr"] & 0 & \\
& D(R) & & & M\otimes\mathbb Z_{\ell\,\underline S} & & .
\end{tikzcd}
\end{equation}

\noindent
印本中 (\SGARefLink{sga7.sga7i.eix.n12-8-9}{12.8.9}) 上方的等同漏去了 $\underline{\operatorname{Hom}}(M,T_\ell(\mathbb G_{m\underline S}))$ 后最外层的右括号; 这里已补齐括号. 注意, $R$ 典范自对偶, 这由自对偶 (\SGARefLink{sga7.sga7i.eix.n12-8-7}{12.8.7}) 及上述正交性断言得出. 与所引构造的唯一区别在于, 所得 $M_\ell\times M_\ell$ 上的配对取值于群
\[
 \bigl(\mathbb G_m(U)/\mathbb G_m(\underline S)\bigr)
 \otimes_{\mathbb Z}\mathbb Z_\ell,
\]
它不再典范同构于 $\mathbb Z_\ell$, 而是通过同构

\begin{equation}\tag{\SGARefTarget{sga7.sga7i.eix.n12-8-10}{12.8.10}}
 \mathbb Z^L\simeq
 \mathbb G_m(U)/\mathbb G_m(\underline S).
\end{equation}

\noindent
典范同构于 $\mathbb Z_\ell^{\,L}$. 此同构将 $\mathbb Z^L$ 的典范基中以 $\lambda$ 为指标的元素映到右端 $T_{i_\lambda}$ 的类, 其中 $T_{i_\lambda}$ 为除子 $D_\lambda$ 的方程. 因而得到典范配对

\begin{equation}\tag{\SGARefTarget{sga7.sga7i.eix.n12-8-11}{12.8.11}}
 \underline{u}_\ell:M_\ell\otimes M_\ell
 \longrightarrow\mathbb Z_\ell^{\,L}\qquad .
\end{equation}

\noindent
几乎显然, 要研究的 $S$ 上单值配对 $u_\ell$ (\SGARefLink{sga7.sga7i.eix.n12-4-7}{12.4.7}) 是 $\underline{u}_\ell$ 与同态 $\mathbb Z_\ell^{\,L}\longrightarrow\mathbb Z_\ell$ 的复合, 后一个同态由下列同态在 $\mathbb Z$ 上与 $\mathbb Z_\ell$ 作张量积得到:

\begin{center}
\begin{tikzcd}
 \mathbb G_m(U)/\mathbb G_m(\underline S)
   \arrow[r] \arrow[d,dash,"\wr"'] &
 \mathbb G_m(\eta)/\mathbb G_m(S)
   \arrow[d,dash,"\wr"'] \\
 \mathbb Z^L & \mathbb Z
\end{tikzcd}
\end{center}

\noindent
这里的同态由 $g:S\longrightarrow\underline S$ 诱导 (它将 $\eta$ 映入 $U$). 这一同态由元素

\begin{equation}\tag{\SGARefTarget{sga7.sga7i.eix.n12-8-12}{12.8.12}}
 (d_\lambda)_{\lambda\in L}\in\mathbb Z^L,\qquad
 d_\lambda\geqslant1\quad\text{对每个 }\lambda\in L,
\end{equation}

\noindent
给出, 其中对每个 $\lambda\in L$, $d_\lambda$ 是 $\underline S$ 上的除子 $D_\lambda$ 在 $S$ 中的逆像的重数. 实际上, 它仅依赖于作为 $V$-代数的完备化 $\widehat{\underline{\mathcal O}}_{X,x_\lambda}$ (第 VI 篇).\footnote{印本在本句中两次引用第 VI 篇, 具体节定位都留空.} 此外, 已知 (第 VI 篇), 对给定的 $\lambda$, $d_\lambda=1$ 当且仅当 $X$ 在 $x_\lambda$ 处\underline{正则}.\footnote{印本中“正则”的主语为小写字母 $x$. 下一句及数学陈述要求这里为概形 $X$.} 由于 \SGARefLink{sga7.sga7i.eix.n12-5}{12.5} 假设 $X$ 正则, 即各 $d_\lambda$ 均等于 $1$, 断言 \SGARefLink{sga7.sga7i.eix.n12-5}{12.5} 将由下列引理得到, 其中已不再出现 $X,S,g$.

\medskip
\noindent
\underline{引理} \SGARefTarget{sga7.sga7i.eix.n12-8-13}{12.8.13}. \underline{若 $\lambda\in L$, 则 (\SGARefLink{sga7.sga7i.eix.n12-8-11}{12.8.11}) 与 $\operatorname{pr}_\lambda$ 的复合等于 (\SGARefLink{sga7.sga7i.eix.n12-4-4}{12.4.4}) 中的型 $\varphi_{x_\lambda}\otimes\varphi_{x_\lambda}$.}

\medskip
为证明最后这个断言, 考虑 $\underline S$ 在 $D_\lambda$ 的一般点处的严格局部化的完备化 $\underline S^\lambda$,\footnote{在此归约中, 印本漏去了 $D_\lambda$ 的下标 $\lambda$, 又将下述纤维积的右因子 $X$ 印成不带下划线. 这里恢复这两个标记: 局部化位于分支 $D_\lambda$ 处, 纤维积由 $\underline S$ 上的泛曲线 $\underline X$ 构成.} 并置
\[
 \underline X^\lambda=
 \underline X\times_{\underline S}\underline S^\lambda.
\]
此时特殊纤维 $\underline X_o^\lambda$ 仅有一个奇点 $y_\lambda$, 它仍为普通二次奇点. 在相应的单值配对

\begin{equation}\tag{\SGARefTarget{sga7.sga7i.eix.n12-8-14}{12.8.14}}
 u_\ell^\lambda:M_\ell^\lambda\otimes M_\ell^\lambda
 \longrightarrow\mathbb Z_\ell\qquad ,
\end{equation}

\noindent
中, $M^\lambda$ 因而秩为 $1$ 或 $0$. 以下容易验证的事实留给读者证明; 它们在我们所需的特例下, 显式给出了此前各构造在局部化下的行为.

\medskip
\noindent
a) $P\times_{\underline S}D_\lambda$ 有唯一的子环面 $T_\lambda$, 使其限制到 $D_\lambda$ 的一般点 $s_\lambda$ 后, 为 $P_{s_\lambda}$ 的极大环面 (维数为 $0$ 或 $1$).

\medskip
\noindent
b) 考虑包含 $T_{\lambda o}\hookrightarrow T_o$ 在特征标群上诱导的转置同态 $M\longrightarrow M^\lambda$, 其中 $T_o$ 为 $P_o$ 的极大环面, 再考虑如 \SGARefLink{sga7.sga7i.eix.n12-4-3}{12.4.3} 那样在相差符号的意义下定义的型 $\varphi_{y_\lambda}:M^\lambda\longrightarrow\mathbb Z$ (将 $X$ 换成 $\underline X^\lambda$). 则在相差符号的意义下, 下图交换:

\begin{equation}\tag{\SGARefTarget{sga7.sga7i.eix.n12-8-15}{12.8.15}}
\begin{tikzcd}[column sep=huge,row sep=large]
 M \arrow[rr] \arrow[dr,"\varphi_{x_\lambda}"']
 & & M^\lambda \arrow[dl,"\varphi_{y_\lambda}"] \\
 & \mathbb Z & \qquad .
\end{tikzcd}
\end{equation}

\noindent
工作转录稿将 (\SGARefLink{sga7.sga7i.eix.n12-8-15}{12.8.15}) 中右侧斜箭头的标号误读为 $\varphi_{\lambda_\lambda}$. 印本图、紧前的定义及下方的公式均为 $\varphi_{y_\lambda}$. (事实上, 在 $X_o$ 中经过 $x_\lambda$ 的两条分枝上选定次序, 会决定 $\underline X_o^\lambda$ 的点 $y_\lambda$ 处的类似选取; 对这样相对应的选取, 上图交换.)

\medskip
\noindent
c) 下图交换:

\begin{center}
\begin{tikzcd}[column sep=huge,row sep=large]
 M_\ell\otimes M_\ell
   \arrow[rr] \arrow[dr,"\operatorname{pr}_\lambda\underline{u}_\ell"']
 & & M_\ell^\lambda\otimes M_\ell^\lambda
   \arrow[dl,"u_\ell^\lambda"] \\
 & \mathbb Z_\ell & \qquad ,
\end{tikzcd}
\end{center}

\noindent
其中记号与 \SGARefLink{sga7.sga7i.eix.n12-8-13}{12.8.13} 和 (\SGARefLink{sga7.sga7i.eix.n12-8-14}{12.8.14}) 相同. 这些关系立即表明, 为证明 \SGARefLink{sga7.sga7i.eix.n12-8-13}{12.8.13}, 只需证明 $\underline X^\lambda$ 的类似关系:
\[
 u_\ell^\lambda=
 (\varphi_{y_\lambda}\otimes\varphi_{y_\lambda})
 \otimes_{\mathbb Z}\mathbb Z_\ell,
\]
这就是对 $\underline S^\lambda$ 上的 $\underline X^\lambda$ 的断言 \SGARefLink{sga7.sga7i.eix.n12-5}{12.5}. 而 $\underline S^\lambda$ 为\underline{零特征}性状.\footnote{在这最后一步中, 印本写成 $\underline S$, 而非整个归约中所用的局部化性状 $\underline S^\lambda$. 英文版恢复了数学上所需的上标.} 上述归约至此完成.

\medskip
\noindent
\underline{注 \SGARefTarget{sga7.sga7i.eix.n12-9}{12.9}.} 注意, 这一归约同时将我们归结为 $X_o$ 仅有一个奇点的情形, 即 Picard--Lefschetz 公式经典上考虑的情形. 还应注意, 此时 $X_o$ 有一个或两个不可约分支, 分别对应于 $\operatorname{rank}M=1$ 和 $\operatorname{rank}M=0$ (后一情形即泛纤维的 Jacobian 有良约化的情形).

\medskip
\noindent
\SGARefTarget{sga7.sga7i.eix.n12-10}{12.10}. \SGARefLink{sga7.sga7i.eix.n12-8}{12.8} 所给的证明可得一个比 \SGARefLink{sga7.sga7i.eix.n12-5}{12.5} 更一般的结果, 无需假设 $X$ 正则: 引入 (\SGARefLink{sga7.sga7i.eix.n12-8-12}{12.8.12}) 中的整数 $d_\lambda\geqslant1$, 此时所有单值配对都由下列配对导出:

\begin{equation}\tag{\SGARefTarget{sga7.sga7i.eix.n12-10-1}{12.10.1}}
 u=\sum_{\lambda\in L}d_\lambda\,
   \varphi_{x_\lambda}\otimes\varphi_{x_\lambda}:
   M\otimes M\longrightarrow\mathbb Z\qquad .
\end{equation}

% END INCLUDED SGA7I_English_current_buildable_source_20260823/components/177_expose_IX_section_12_8_9_through_12_10.tex

